% Generated by scripts/assemble.py; corrections belong in source/corrections.json.
\examhead{74}{东北大学 2026 年数学分析试卷}{73}
% p074-b11; source PDF page 74
\blockanchor{p074-b11}
\question{1}(15 分) 求极限 $\lim_{n\to \infty}\left[\ln \left(1 + \frac{1}{n}\right)^n\right]^n.$\par
% p074-b12; source PDF page 74
\blockanchor{p074-b12}
\question{2}(15 分) 将 $f(x) = \ln x$ 按 $\frac{x - 1}{x + 1}$ 进行幂级数展开.\par
% p074-b13; source PDF page 74
\blockanchor{p074-b13}
\question{3}(15 分) 设 $m, n \in \mathbb{N}^{+}$, 求 $I(m, n) = \int_{0}^{1} (\ln x)^{m} x^{n} \, \mathrm{d}x$.\par
% p074-b14; source PDF page 74
\blockanchor{p074-b14}
\question{4}(15 分) 计算三重积分 $\iiint_{\Omega} e^{|z|} \, \mathrm{d}x \, \mathrm{d}y \, \mathrm{d}z$, 其中 $\Omega$ 是单位球.\par
% p074-b15; source PDF page 74
\blockanchor{p074-b15}
\question{5}(15 分) 叙述定义: $x \to +\infty$ 时 $f(x) \to -\infty$; $x \to 0^{+}$ 时, $f(x)$ 不以 $A$ 为极限.\par
% p074-b16; source PDF page 74
\blockanchor{p074-b16}
\question{6}(15 分) 证明级数 $\sum_{n=1}^{\infty}[nxe^{-nx} - (n-1)xe^{-(n-1)x}]$ 在 $[0,1]$ 上收敛但不一致收敛, 而和函数连续.\par
% p074-b17; source PDF page 74
\par\Needspace{6\baselineskip}
\blockanchor{p074-b17}
\question{7}(15 分) 设 $f(x)$ 在 $[x_1, x_2]$ 上连续 $(x_1 x_2 > 0)$, 在 $(x_1, x_2)$ 上可导, 证明: 存在 $\xi \in (x_1, x_2)$, 使得\par
\nopagebreak[4]
% p074-b18; source PDF page 74
\blockanchor{p074-b18}
\displayformula{\frac {1}{x _ {2} - x _ {1}} \left| \begin{array}{l l} x _ {1} & f (x _ {1}) \\ x _ {2} & f (x _ {2}) \end{array} \right| = \xi f ^ {\prime} (\xi) - f (\xi).}
% p074-b19; source PDF page 74
\blockanchor{p074-b19}
\question{8}(15 分) 设 $f(x)$ 在 $(- \infty, + \infty)$ 上二阶连续可微, 且 $f(x) + f''(x) = -xg(x)f'(x)$, 其中 $g(x)$ 为非负函数. 证明: $f(x)$ 有界.\par
% p074-b20; source PDF page 74
\par\Needspace{6\baselineskip}
\blockanchor{p074-b20}
\question{9}(15 分) 设 $f(x)$ 为 $[a, b]$ 上的二阶可导函数, 且 $f''(x) \leq 0$, $f(x) \geq 0$, 证明:\par
\nopagebreak[4]
% p074-b21; source PDF page 74
\blockanchor{p074-b21}
\displayformula{\max _ {x \in [ a, b ]} f (x) \leq \frac {2}{b - a} \int_ {a} ^ {b} f (x) \mathrm{d} x.}
% p074-b22; source PDF page 74
\blockanchor{p074-b22}
\question{10}(15 分) 证明: $f(x) = \begin{cases} \frac{1}{n}, & x = \frac{m}{n} (m, n \text{ 是互质正整数}) \\ 0, & x \text{ 为无理数}. \end{cases}$ 在无理点连续, 在有理点不连续.\par
